\documentclass[10pt,a4paper]{amsart}
\usepackage[margin=19mm]{geometry}
\usepackage{amsmath,amssymb,mathtools}
\usepackage[scr=rsfs,cal=euler]{mathalfa}
\usepackage{xfrac}
\usepackage[unicode,hidelinks,bookmarks=false]{hyperref}
\newcommand{\ie}{i.e.\ }
\newcommand{\cf}{cf.\ }
\newcommand{\resp}{respectively\ }
\newcommand{\Subsection}{\subsection}
\providecommand{\nobreakdash}{\nobreak\hbox{-}}
\setlength{\emergencystretch}{2em}
\multlinegap0pt
\usepackage{xcolor}
\usepackage{amssymb}
\usepackage{bm}
\usepackage{algorithm}\usepackage{algpseudocode}
\usepackage{tikz}
\usepackage{mathtools}
\newtheorem{lemma}{Lemma}
\newcommand{\A}{\mathcal{A}}
\newcommand{\C}{\mathcal{C}}
\newcommand{\R}{\mathcal{R}}
\title{Beyond Hand-Derived Inequalities: Decision Diagrams for Cut Generation in Binary Polynomial Optimization}
\author{Martin Cooper, Margarita Castro}
\date{Source: arXiv:2607.28511v1 (CC BY 4.0)}
\begin{document}
\maketitle

\noindent\emph{Source and licence.} Beyond Hand-Derived Inequalities: Decision Diagrams for Cut Generation in Binary Polynomial Optimization. Version arXiv:2607.28511v1 (CC BY 4.0), distributed by arXiv under the Creative Commons Attribution 4.0 International licence, http://creativecommons.org/licenses/by/4.0/, as recorded in the arXiv licence metadata for this exact version and shown on the abstract page https://arxiv.org/abs/2607.28511v1. Full licence notice: \texttt{licenses/arxiv-2607.28511-CC-BY-4.0.txt}.\par\smallskip

\noindent\textit{Editorial scope.} Counted unit: the article's lemma lem:valid\_antichains (source0.tex lines 430--433). The article's complete original proof of it is delivered with it (source0.tex lines 435--456, one \texttt{proof} environment of 22 lines). No mathematics is written, repaired or re-derived: the statement and the proof are the article's own lines, in the article's own notation, and no retained line is substituted or re-flowed.  No further source-local context is needed: every cross-reference of every retained line resolves inside the two delivered spans.  Nothing was omitted from any retained span, so the package carries no omission record.  No retained line contains a citation, so the wrapper carries no bibliography and no bibliography entry.  No retained line contains a figure environment, an image-inclusion command or drawing code, so no image is delivered and the package declares no asset dependency.  Editorial formatting only, in addition: an AMS \texttt{amsart} preamble that restates the article's own declarations and macros used by the retained lines, namely the environments \texttt{lemma} on separate counters, together with the standard packages those lines need.  The retained environments print in this document as Lemma 1 in order of appearance, whereas the article prints the same items with the numbering of its own full document.  The complete native source of the article as distributed in the arXiv e-print for this exact version is bundled unchanged as \texttt{sources/206424/source0.tex}.
\begin{lemma}
\label{lem:valid_antichains}
Consider a vertex order $\sigma$ and a layer $N_i$ of the compact DD $B=(N,A)$. The reachable states after assigning $U_i$ variables are in bijection with the valid antichains of $(\R_i,\subseteq)$.
\end{lemma}

\begin{proof}
We first show that each reachable state in $N_i$ corresponds to a valid antichain in $(\R_i,\subseteq)$. Let $s_i$ be a reachable state and define $I(s_i)=\{\rho_i(e):e\in s_i\}\subseteq\R_i$ as the set of reduced active hyperedges that are compatible in a reachable state $s_i$. Set $I(s_i)$ is downward closed because if $R'\in I(s_i)$ and $R\in\R_i$ satisfies $R\subseteq R'$, then all vertices of $R$ also have value one under the current assignment, so $R$ is compatible as well. Therefore $R\in I(s_i)$, and $I(s_i)$ is an ideal.

Since every ideal is determined by its maximal elements, let us define
$
\operatorname{Max}(I(s_i)) = \{R\in I(s_i): \nexists R'\in I(s_i)\text{ with }R\subsetneq R'\}
$
as the set of maximal elements, which in turn form an antichain. Note that the elements in $\operatorname{Max}(I(s_i))$ are also valid: if
$R\in\R_i$ satisfies $R\subseteq X(\operatorname{Max}(I(s_i)))$, then all
vertices of $R$ have value one, so $R$ is compatible and hence belongs to $I(s_i)$. Because $I(s_i)$ is an ideal, $R$ is contained in some maximal element of $I(s_i)$.

Now, we show the other direction by assuming that $\C$ is a valid antichain. Consider the assignment on $U_i$ that sets the vertices in $X(\C)$ to one and the remaining vertices of $U_i$ to zero. Under this assignment, a reduced active hyperedge $R\in\R_i$ is
compatible if and only if $R\subseteq X(\C)$. By validity, every such $R$ is
contained in some member of $\C$. Hence, the induced state is
$
s_i(\C)=
\{e\in\A_i :
\rho_i(e)\subseteq M \text{ for some } M\in\C\}.
$
The maximal elements of the corresponding ideal are exactly the members of
$\C$. This gives the claimed bijection. \hfill$\square$
\end{proof}

\end{document}
